%
\begin{isabellebody}%
\setisabellecontext{Week{\isacharunderscore}{\kern0pt}{\isadigit{2}}}%
%
\isadelimtheory
%
\endisadelimtheory
%
\isatagtheory
%
\endisatagtheory
{\isafoldtheory}%
%
\isadelimtheory
%
\endisadelimtheory
%
\begin{isamarkuptext}%
\ExerciseSheet{2}{}%
\end{isamarkuptext}\isamarkuptrue%
%
\begin{isamarkuptext}%
Important Note: Please download Isabelle from the link provided in the slides and bring a 
      laptop to the large group tutorial. Performing the proofs in Isabelle will be an integral part
      of those tutorials.%
\end{isamarkuptext}\isamarkuptrue%
%
\begin{isamarkuptext}%
Before you start, please open a new Isabelle/HOL theory file called Week\_1.thy and add the
       following to its beginning.

\noindent theory Week\_2\\
imports Main\\
begin%
\end{isamarkuptext}\isamarkuptrue%
%
\begin{isamarkuptext}%
\Exercise{Addition is Associative}%
\end{isamarkuptext}\isamarkuptrue%
%
\begin{isamarkuptext}%
Consider the function \isa{{\isacharparenleft}{\kern0pt}{\isacharplus}{\kern0pt}{\isacharparenright}{\kern0pt}} defined in Isabelle/HOL for natural numbers.

Prove the following theorem, first pen-and-paper, and then formally in Isabelle.%
\end{isamarkuptext}\isamarkuptrue%
\isakeywordONE{lemma}\isamarkupfalse%
\ add{\isacharunderscore}{\kern0pt}associative{\isacharcolon}{\kern0pt}\ {\isachardoublequoteopen}{\isacharparenleft}{\kern0pt}{\isacharparenleft}{\kern0pt}a{\isacharcolon}{\kern0pt}{\isacharcolon}{\kern0pt}nat{\isacharparenright}{\kern0pt}\ {\isacharplus}{\kern0pt}\ b{\isacharparenright}{\kern0pt}\ {\isacharplus}{\kern0pt}\ c\ {\isacharequal}{\kern0pt}\ a\ {\isacharplus}{\kern0pt}\ {\isacharparenleft}{\kern0pt}b\ {\isacharplus}{\kern0pt}\ c{\isacharparenright}{\kern0pt}{\isachardoublequoteclose}%
\begin{isamarkuptext}%
\paragraph{Hints:}

 1. Perform the proof by induction, using substitutions \isa{subst}, backward reasoning using the method
   \isa{rule}, and forward reasoning using \isa{OF}.

 2. Think about which variable to perform induction on. Addition on natural numbers is defined by
    recursion on its first argument: look at \isa{thm\ add{\isacharunderscore}{\kern0pt}{\isadigit{0}}{\isacharunderscore}{\kern0pt}left\ add{\isacharunderscore}{\kern0pt}Suc}.

 3. You will need to identify multiple lemmas about how \isa{add} works. Do not prove those lemmas. Use
    \isa{sorry} as the proof. This is the top-down approach.

 4. There is a proof which uses only the following lemmas: \isa{add{\isacharunderscore}{\kern0pt}{\isadigit{0}}{\isacharunderscore}{\kern0pt}left{\isacharcomma}{\kern0pt}\ add{\isacharunderscore}{\kern0pt}Suc{\isacharcomma}{\kern0pt}\ nat{\isachardot}{\kern0pt}inject{\isacharcomma}{\kern0pt}\ iffD{\isadigit{2}}}%
\end{isamarkuptext}\isamarkuptrue%
%
\isadelimproof
%
\endisadelimproof
%
\isatagproof
%
\endisatagproof
{\isafoldproof}%
%
\isadelimproof
%
\endisadelimproof
%
\begin{isamarkuptext}%
\Exercise{Multiplication is Monotone}%
\end{isamarkuptext}\isamarkuptrue%
%
\begin{isamarkuptext}%
Consider the multiplication function \isa{{\isacharparenleft}{\kern0pt}{\isacharasterisk}{\kern0pt}{\isacharparenright}{\kern0pt}} defined in Isabelle. Prove the following
      property of it, i.e.\ that it is monotonically increasing in both arguments. You should prove
      it both, pen-and-paper and in Isabelle.%
\end{isamarkuptext}\isamarkuptrue%
\isakeywordONE{lemma}\isamarkupfalse%
\ mult{\isacharunderscore}{\kern0pt}mono{\isacharcolon}{\kern0pt}\isanewline
\ \ \isakeywordTWO{fixes}\ a\ b\ c\ d{\isacharcolon}{\kern0pt}{\isacharcolon}{\kern0pt}\ nat\ %
\isamarkupcmt{Note the alternative way of fixing the types of the constants%
}\isanewline
\ \ \isakeywordTWO{assumes}\ {\isachardoublequoteopen}a\ {\isasymle}\ b{\isachardoublequoteclose}\ {\isachardoublequoteopen}c\ {\isasymle}\ d{\isachardoublequoteclose}\isanewline
\ \ \isakeywordTWO{shows}\ {\isachardoublequoteopen}a\ {\isacharasterisk}{\kern0pt}\ c\ {\isasymle}\ b\ {\isacharasterisk}{\kern0pt}\ d{\isachardoublequoteclose}%
\begin{isamarkuptext}%
Your pen-and-paper proof should indicate

 1. what lemmas are you assuming,

 2. on what variable are you preforming induction,

 3. what are assumptions and the proof goals in the base case and the step case, and what is the
    induction hypothesis, and

 4. how are you instantiating the induction hypothesis, i.e. how do you show that its assumptions
    hold and which quantified variables are instantiated with which constants.

\paragraph{Hints:}
 1. The proof should be by induction

 2. The following lemmas suffice to prove the goal with only substitution, forward reasoning using
     \isa{OF}, backward reasoning using \isa{rule}, and induction:

     \isa{le{\isacharunderscore}{\kern0pt}{\isadigit{0}}{\isacharunderscore}{\kern0pt}eq{\isacharcomma}{\kern0pt}\ iffD{\isadigit{1}}{\isacharcomma}{\kern0pt}\ mult{\isacharunderscore}{\kern0pt}zero{\isacharunderscore}{\kern0pt}right{\isacharcomma}{\kern0pt}\ le{\isacharunderscore}{\kern0pt}refl{\isacharcomma}{\kern0pt}\ not{\isacharunderscore}{\kern0pt}le{\isacharcomma}{\kern0pt}\ Suc{\isacharunderscore}{\kern0pt}leI{\isacharcomma}{\kern0pt}\ le{\isacharunderscore}{\kern0pt}antisym{\isacharcomma}{\kern0pt}\ le{\isacharunderscore}{\kern0pt}SucI{\isacharcomma}{\kern0pt}\ mult{\isacharunderscore}{\kern0pt}le{\isacharunderscore}{\kern0pt}mono{\isadigit{1}}{\isacharcomma}{\kern0pt}\ mult{\isacharunderscore}{\kern0pt}le{\isacharunderscore}{\kern0pt}mono{\isadigit{2}}{\isacharcomma}{\kern0pt}\ le{\isacharunderscore}{\kern0pt}trans}

 3. In the step case, perform a case analysis with \isa{proof\ {\isacharparenleft}{\kern0pt}cases\ {\isachardoublequote}{\kern0pt}c\ {\isasymle}\ d{\isachardoublequote}{\kern0pt}{\isacharparenright}{\kern0pt}}. In the second case,
     together with the assumption \isa{c\ {\isasymle}\ Suc\ d}, you can derive \isa{c\ {\isacharequal}{\kern0pt}\ Suc\ d}.%
\end{isamarkuptext}\isamarkuptrue%
%
\isadelimproof
%
\endisadelimproof
%
\isatagproof
%
\endisatagproof
{\isafoldproof}%
%
\isadelimproof
%
\endisadelimproof
%
\begin{isamarkuptext}%
\Exercise{Induction Hypothesis Generalisation}%
\end{isamarkuptext}\isamarkuptrue%
%
\begin{isamarkuptext}%
Consider the following function, which adds two natural numbers using an accumulator: in
      each recursive call it moves one \isa{\isaconst{Suc}} from its first argument to its second.%
\end{isamarkuptext}\isamarkuptrue%
\isakeywordONE{fun}\isamarkupfalse%
\ add{\isacharunderscore}{\kern0pt}acc\ {\isacharcolon}{\kern0pt}{\isacharcolon}{\kern0pt}\ {\isachardoublequoteopen}nat\ {\isasymRightarrow}\ nat\ {\isasymRightarrow}\ nat{\isachardoublequoteclose}\ \isakeywordTWO{where}\isanewline
\ \ {\isachardoublequoteopen}add{\isacharunderscore}{\kern0pt}acc\ {\isadigit{0}}\ m\ {\isacharequal}{\kern0pt}\ m{\isachardoublequoteclose}\isanewline
{\isacharbar}{\kern0pt}\ {\isachardoublequoteopen}add{\isacharunderscore}{\kern0pt}acc\ {\isacharparenleft}{\kern0pt}Suc\ n{\isacharparenright}{\kern0pt}\ m\ {\isacharequal}{\kern0pt}\ add{\isacharunderscore}{\kern0pt}acc\ n\ {\isacharparenleft}{\kern0pt}Suc\ m{\isacharparenright}{\kern0pt}{\isachardoublequoteclose}%
\begin{isamarkuptext}%
We want to prove that \isa{add{\isacharunderscore}{\kern0pt}acc} computes addition.%
\end{isamarkuptext}\isamarkuptrue%
\isakeywordONE{lemma}\isamarkupfalse%
\ add{\isacharunderscore}{\kern0pt}acc{\isacharunderscore}{\kern0pt}correct{\isacharcolon}{\kern0pt}\ {\isachardoublequoteopen}add{\isacharunderscore}{\kern0pt}acc\ n\ m\ {\isacharequal}{\kern0pt}\ n\ {\isacharplus}{\kern0pt}\ m{\isachardoublequoteclose}%
\begin{isamarkuptext}%
\paragraph{(a)} First try to prove the lemma by plain induction on \isa{n}, i.e.\ with
  \isa{proof{\isacharparenleft}{\kern0pt}induction\ n{\isacharparenright}{\kern0pt}}. Write down the goal of the step case and the induction hypothesis. Explain,
  pen-and-paper, why the induction hypothesis is not strong enough to prove the step case.

 \paragraph{(b)} Now prove the lemma, first pen-and-paper and then formally in Isabelle, by
  generalising the induction hypothesis over \isa{m}, i.e.\ with \isa{proof{\isacharparenleft}{\kern0pt}induction\ n\ arbitrary{\isacharcolon}{\kern0pt}\ m{\isacharparenright}{\kern0pt}}.
  Your pen-and-paper proof should state the generalised induction hypothesis and say with which term
  you instantiate its universally quantified variable in the step case.

\paragraph{Hints:}
 1. \isa{add{\isacharunderscore}{\kern0pt}acc} is a recursive function defined with \isa{fun}, like \isa{add} in Section 2.2.2 of
    Concrete Semantics. Its defining equations are available as \isa{add{\isacharunderscore}{\kern0pt}acc{\isachardot}{\kern0pt}simps{\isacharparenleft}{\kern0pt}{\isadigit{1}}{\isacharparenright}{\kern0pt}} and
    \isa{add{\isacharunderscore}{\kern0pt}acc{\isachardot}{\kern0pt}simps{\isacharparenleft}{\kern0pt}{\isadigit{2}}{\isacharparenright}{\kern0pt}} (look at \isa{thm\ add{\isacharunderscore}{\kern0pt}acc{\isachardot}{\kern0pt}simps}), and can be used with \isa{subst} like any
    other equation.

 2. Use \isa{thm\ Suc{\isachardot}{\kern0pt}IH} in the step case to compare the induction hypothesis you get with and without
    \isa{arbitrary{\isacharcolon}{\kern0pt}\ m}.

 3. An instance of a universally quantified fact can be obtained with \isa{of}, e.g.\
    \isa{Suc{\isachardot}{\kern0pt}IH{\isacharbrackleft}{\kern0pt}of\ {\isachardoublequote}{\kern0pt}Suc\ m{\isachardoublequote}{\kern0pt}{\isacharbrackright}{\kern0pt}}.

 4. The following lemmas suffice: \isa{add{\isacharunderscore}{\kern0pt}acc{\isachardot}{\kern0pt}simps{\isacharcomma}{\kern0pt}\ add{\isacharunderscore}{\kern0pt}{\isadigit{0}}{\isacharunderscore}{\kern0pt}left{\isacharcomma}{\kern0pt}\ add{\isacharunderscore}{\kern0pt}Suc{\isacharcomma}{\kern0pt}\ add{\isacharunderscore}{\kern0pt}Suc{\isacharunderscore}{\kern0pt}right}.%
\end{isamarkuptext}\isamarkuptrue%
%
\isadelimproof
%
\endisadelimproof
%
\isatagproof
%
\endisatagproof
{\isafoldproof}%
%
\isadelimproof
%
\endisadelimproof
%
\begin{isamarkuptext}%
\paragraph{(c) Challenge (optional)} The following function multiplies using an accumulator.
  Prove, pen-and-paper and in Isabelle, that it computes \isa{acc\ {\isacharplus}{\kern0pt}\ n\ {\isacharasterisk}{\kern0pt}\ m}. Which variable must you
  generalise over this time, and why?%
\end{isamarkuptext}\isamarkuptrue%
\isakeywordONE{fun}\isamarkupfalse%
\ mult{\isacharunderscore}{\kern0pt}acc\ {\isacharcolon}{\kern0pt}{\isacharcolon}{\kern0pt}\ {\isachardoublequoteopen}nat\ {\isasymRightarrow}\ nat\ {\isasymRightarrow}\ nat\ {\isasymRightarrow}\ nat{\isachardoublequoteclose}\ \isakeywordTWO{where}\isanewline
\ \ {\isachardoublequoteopen}mult{\isacharunderscore}{\kern0pt}acc\ {\isadigit{0}}\ m\ acc\ {\isacharequal}{\kern0pt}\ acc{\isachardoublequoteclose}\isanewline
{\isacharbar}{\kern0pt}\ {\isachardoublequoteopen}mult{\isacharunderscore}{\kern0pt}acc\ {\isacharparenleft}{\kern0pt}Suc\ n{\isacharparenright}{\kern0pt}\ m\ acc\ {\isacharequal}{\kern0pt}\ mult{\isacharunderscore}{\kern0pt}acc\ n\ m\ {\isacharparenleft}{\kern0pt}acc\ {\isacharplus}{\kern0pt}\ m{\isacharparenright}{\kern0pt}{\isachardoublequoteclose}\isanewline
\isanewline
\isakeywordONE{lemma}\isamarkupfalse%
\ mult{\isacharunderscore}{\kern0pt}acc{\isacharunderscore}{\kern0pt}correct{\isacharcolon}{\kern0pt}\ {\isachardoublequoteopen}mult{\isacharunderscore}{\kern0pt}acc\ n\ m\ acc\ {\isacharequal}{\kern0pt}\ acc\ {\isacharplus}{\kern0pt}\ n\ {\isacharasterisk}{\kern0pt}\ m{\isachardoublequoteclose}%
\begin{isamarkuptext}%
\paragraph{Hint:} The following lemmas suffice:
      \isa{mult{\isacharunderscore}{\kern0pt}acc{\isachardot}{\kern0pt}simps{\isacharcomma}{\kern0pt}\ mult{\isacharunderscore}{\kern0pt}{\isadigit{0}}{\isacharcomma}{\kern0pt}\ add{\isacharunderscore}{\kern0pt}{\isadigit{0}}{\isacharunderscore}{\kern0pt}right{\isacharcomma}{\kern0pt}\ mult{\isacharunderscore}{\kern0pt}Suc{\isacharcomma}{\kern0pt}\ add{\isachardot}{\kern0pt}assoc}.%
\end{isamarkuptext}\isamarkuptrue%
%
\isadelimproof
%
\endisadelimproof
%
\isatagproof
\isanewline
%
\endisatagproof
{\isafoldproof}%
%
\isadelimproof
%
\endisadelimproof
%
\isadelimtheory
%
\endisadelimtheory
%
\isatagtheory
%
\endisatagtheory
{\isafoldtheory}%
%
\isadelimtheory
%
\endisadelimtheory
%
\end{isabellebody}%
\endinput
%:%file=/media/llms/Theorem_Proving/7CCSACTP-private/Week_2/ex/Week_2.thy%:%
%:%19=6%:%
%:%23=17%:%
%:%24=18%:%
%:%25=19%:%
%:%29=22%:%
%:%30=23%:%
%:%31=24%:%
%:%32=25%:%
%:%33=26%:%
%:%34=27%:%
%:%38=30%:%
%:%42=34%:%
%:%43=35%:%
%:%44=36%:%
%:%46=49%:%
%:%47=49%:%
%:%49=51%:%
%:%50=52%:%
%:%51=53%:%
%:%52=54%:%
%:%53=55%:%
%:%54=56%:%
%:%55=57%:%
%:%56=58%:%
%:%57=59%:%
%:%58=60%:%
%:%59=61%:%
%:%60=62%:%
%:%77=111%:%
%:%81=113%:%
%:%82=114%:%
%:%83=115%:%
%:%85=140%:%
%:%86=140%:%
%:%87=141%:%
%:%88=141%:%
%:%89=141%:%
%:%90=142%:%
%:%91=143%:%
%:%93=146%:%
%:%94=147%:%
%:%95=148%:%
%:%96=149%:%
%:%97=150%:%
%:%98=151%:%
%:%99=152%:%
%:%100=153%:%
%:%101=154%:%
%:%102=155%:%
%:%103=156%:%
%:%104=157%:%
%:%105=158%:%
%:%106=159%:%
%:%107=160%:%
%:%108=161%:%
%:%109=162%:%
%:%110=163%:%
%:%111=164%:%
%:%111=165%:%
%:%112=166%:%
%:%113=167%:%
%:%114=168%:%
%:%131=225%:%
%:%135=227%:%
%:%136=228%:%
%:%138=230%:%
%:%139=230%:%
%:%140=231%:%
%:%141=232%:%
%:%143=234%:%
%:%145=236%:%
%:%146=236%:%
%:%148=239%:%
%:%149=240%:%
%:%150=241%:%
%:%151=242%:%
%:%152=243%:%
%:%153=244%:%
%:%154=245%:%
%:%155=246%:%
%:%156=247%:%
%:%157=248%:%
%:%158=249%:%
%:%159=250%:%
%:%160=251%:%
%:%161=252%:%
%:%162=253%:%
%:%163=254%:%
%:%164=255%:%
%:%165=256%:%
%:%166=257%:%
%:%167=258%:%
%:%168=259%:%
%:%169=260%:%
%:%186=311%:%
%:%187=312%:%
%:%188=313%:%
%:%190=315%:%
%:%191=315%:%
%:%192=316%:%
%:%193=317%:%
%:%194=318%:%
%:%195=319%:%
%:%196=319%:%
%:%198=321%:%
%:%199=322%:%
%:%207=364%:%
